\begin{lemma}[invariants]
\label{lem:invariants}
After \textsc{Load} succeeds, the following hold at every point of the solve.
\begin{enumerate}
\item[(I1)] A cell not in $U$ has at most one candidate.
\item[(I2)] For each mask, once \textsc{UpdateBand} has run on it, $M_{d,b}\subseteq K_{d,b}$; before that,
  $K_{d,b}=0$.
\item[(I3)] If the bit of row $R$ in $O_d$ is clear, then $d$ has at most one candidate cell in row $R$, and that cell
  has no other candidate.
\end{enumerate}
\end{lemma}
\begin{proof}
(I1) Cells leave $U$ only in \textsc{Place}, in a row-single placement and in a naked-single placement. \textsc{Place}
removes every other digit from the cell, a row-single placement clears the cell from all nine digits and then restores $d$'s mask as it was just before, so
the cell keeps at most $d$, and a naked single already has one candidate. By monotonicity the cell never gains candidates.
(I2) \textsc{UpdateBand} sets $K_{d,b}$ to the current mask; afterwards the mask only shrinks and $K_{d,b}$ changes only
at the next \textsc{UpdateBand} on the same mask. (I3) Row bits of $O_d$ are cleared by the same three placements, each
of which leaves the placed cell $y$ with no candidate other than $d$ and removes $d$ from the rest of $y$'s row: a row
single because the row had no other candidate, \textsc{Place} and a naked-single placement through
\textsc{AfterPlacing}. Monotonicity preserves both properties.
\end{proof}

